\chapter{面向边缘部署需求的YOLO11-StackLite模型压缩方法}{Model Compression Method of YOLO11-StackLite for Edge Deployment}

第2章提出的 YOLO11-StackLite 模型在复杂堆叠货物场景下取得了较好的检测、分割与盘点效果。然而，随着模型精度提升，网络结构中仍然存在一定的参数冗余与计算冗余，这会在一定程度上限制模型在边缘计算场景中的部署应用。对于立体仓库库存盘点任务而言，系统通常要求模型在保证识别精度的同时具备较低的存储开销和较高的推理效率，因此有必要进一步开展模型压缩研究。

模型压缩的目标是在尽量保持任务性能稳定的前提下，降低模型参数量与计算复杂度，从而提升模型的部署友好性。考虑到本文当前研究重点在于提升模型在边缘场景中的应用可行性，而非针对特定硬件平台开展部署优化，因此本章主要从模型参数量、浮点运算量以及检测、分割与盘点相关指标等角度，对模型压缩效果进行间接评估。基于此，本文针对第2章构建的 YOLO11-StackLite，采用结构化剪枝降低模型复杂度，并结合知识蒸馏缓解压缩带来的性能退化，构建适用于堆叠货物盘点任务的轻量化模型。

\section{模型压缩需求分析}{Analysis of Model Compression Requirements}

YOLO11-StackLite 在第2章实验中已表现出较优的检测与分割性能，但其作为实例分割模型，仍包含较多卷积层与特征融合操作。在复杂堆叠货物场景中，模型不仅需要提取目标轮廓信息，还需要保留层间结构信息以支持后续数量推理，因此网络结构通常较普通目标检测模型更为复杂。虽然该结构有助于提升识别精度，但也会带来较高的参数开销与理论计算开销。

在实际仓储盘点任务中，边缘侧设备往往具有算力有限、显存容量受限以及实时处理要求较高等特点。若模型参数规模过大、推理计算量过高，则可能增加部署成本，并影响系统响应效率。因此，有必要在尽量保持检测、分割和盘点性能的前提下，对模型进行轻量化处理。基于这一需求，本章从结构化剪枝和知识蒸馏两个方面开展研究：前者用于降低模型结构复杂度，后者用于弥补压缩过程中可能出现的性能损失。

\section{基于结构化剪枝的轻量化方法}{Lightweight Method Based on Structured Pruning}

\subsection{结构化剪枝原理}{Principle of Structured Pruning}

模型剪枝是深度模型压缩中的常见方法之一，其基本思想是通过删除网络中的冗余结构单元，降低模型规模与计算复杂度。根据剪枝对象的不同，模型剪枝通常可分为非结构化剪枝与结构化剪枝两类。非结构化剪枝主要对单个权重连接进行筛除，虽然能够获得较高的稀疏率，但由于其稀疏结构不规则，实际部署时难以直接转化为显著的硬件加速收益。相比之下，结构化剪枝通常以卷积核、通道或模块为基本单位进行裁剪，压缩后的网络结构更加规则，更有利于后续模型部署与推理优化。

考虑到本文研究目标是提升模型在边缘场景中的部署可行性，故选择结构化剪枝作为主要压缩手段。该方法不仅能够有效降低模型参数量和浮点运算量，而且在工程实现上更便于直接应用于现有推理框架。因此，本文采用基于通道重要性评估的结构化剪枝策略，对 YOLO11-StackLite 中的冗余结构进行压缩。

\subsection{YOLO11-StackLite剪枝对象选择}{Selection of Pruning Targets in YOLO11-StackLite}

在对 YOLO11-StackLite 进行剪枝时，若对整个网络进行无差别压缩，容易导致关键特征表达能力显著下降，进而影响实例分割和后续盘点性能。因此，剪枝对象的选择需要兼顾压缩率与任务性能稳定性。

结合第2章模型结构特点，本文优先将 Neck 中的特征融合层以及部分卷积层作为主要剪枝对象。其原因在于：Neck 负责多尺度特征融合，内部存在一定程度的通道冗余，通过合理剪枝可在不显著破坏基础特征表达能力的前提下降低模型复杂度。对于 Backbone 中承担基础表征任务的关键层，则采用相对保守的剪枝策略，以避免因过度压缩导致底层纹理信息和结构特征损失。此外，对于检测头相关结构，也仅进行适度压缩，以保证输出层的稳定性。

通过上述剪枝对象选择方式，可以在保持模型核心表达能力的同时，尽量释放结构冗余，从而为后续蒸馏精度恢复阶段提供更稳定的学生模型基础。

\subsection{剪枝流程设计}{Pruning Process Design}

本文采用的结构化剪枝流程主要包括以下四个步骤：

\par\noindent
（1）训练原始 YOLO11-StackLite 模型。首先在 TBD 数据集上完成原始模型训练，获得作为压缩基准的高性能模型。

\par\noindent
（2）评估结构单元重要性。根据卷积通道的重要性指标，对网络中待剪枝层的通道进行排序。本文采用基于通道重要性评估的方式识别冗余结构单元，以确定剪枝对象。

\par\noindent
（3）执行结构化剪枝。按照设定的剪枝比例，对选定层中的低重要性通道进行裁剪，得到初始剪枝模型。

\par\noindent
（4）对剪枝模型进行微调。由于结构化剪枝会改变原始模型的特征分布，因此需要对剪枝后的模型进行重新训练或微调，以恢复其基础任务性能。

通过上述过程，本文可获得参数量和计算复杂度均明显下降的轻量化模型。然而，仅采用剪枝通常会带来不同程度的精度下降，因此下一节进一步引入知识蒸馏方法，对压缩模型进行性能补偿。

\section{基于知识蒸馏的精度恢复方法}{Accuracy Recovery Method Based on Knowledge Distillation}

\subsection{知识蒸馏原理}{Principle of Knowledge Distillation}

知识蒸馏的核心思想是利用性能较强的教师模型指导轻量化学生模型训练，使学生模型在较小规模下尽可能继承教师模型的判别能力。在模型压缩任务中，剪枝能够有效降低模型复杂度，但也可能削弱模型的特征表达能力，进而导致检测、分割以及盘点性能下降。为缓解这一问题，本文在剪枝后引入知识蒸馏策略，通过教师模型向学生模型传递输出层响应和中间特征信息，从而提升压缩模型的性能恢复能力。

\subsection{教师模型与学生模型构建}{Construction of Teacher and Student Models}

本文以第2章训练完成的 YOLO11-StackLite 作为教师模型。该模型在 TBD 数据集上具有较好的检测、分割和盘点性能，能够为轻量化模型提供较为充分的知识指导。学生模型则为结构化剪枝后得到的压缩模型，其参数量和计算复杂度明显低于教师模型，但在初始状态下存在一定性能损失。

通过将原始高性能模型作为教师模型、将剪枝模型作为学生模型，可以建立“高性能完整模型 - 轻量化压缩模型”之间的知识迁移关系，从而在保持压缩效果的同时尽可能恢复学生模型的任务性能。

\subsection{蒸馏损失设计}{Design of Distillation Loss}

为提高学生模型对教师模型知识的学习能力，本文在原始任务损失的基础上引入蒸馏损失，构建联合优化目标。原始任务损失主要用于保证学生模型在目标检测与实例分割任务上的基本学习能力；蒸馏损失则用于约束学生模型在输出响应或中间特征层面向教师模型靠近。

设学生模型原始任务损失为 $L_{\mathrm{task}}$，蒸馏损失为 $L_{\mathrm{distill}}$，则总损失可表示为：
\begin{equation}
L = L_{\mathrm{task}} + \lambda L_{\mathrm{distill}}
\end{equation}
其中，$\lambda$ 为蒸馏损失权重系数，用于平衡原始任务学习与知识迁移之间的关系。

在具体实现中，本文可采用以下两类蒸馏信息：
\par\noindent
（1）输出层响应蒸馏。通过约束学生模型与教师模型在最终输出层上的预测响应一致性，增强学生模型的判别能力。

\par\noindent
（2）中间特征蒸馏。通过约束学生模型与教师模型在中间层特征表达上的接近程度，帮助学生模型学习更加稳定的特征表示。

上述蒸馏方式能够在不显著增加推理成本的前提下，提高剪枝模型在复杂堆叠场景下的检测与分割性能，为后续盘点任务提供更可靠的识别基础。

\section{压缩模型构建与训练流程}{Construction and Training Process of the Compressed Model}

综合结构化剪枝与知识蒸馏策略，本文的压缩模型构建流程如下：

\par\noindent
（1）首先在 TBD 数据集上训练完整的 YOLO11-StackLite 模型，作为后续压缩与蒸馏的基准模型；

\par\noindent
（2）然后基于通道重要性评估结果，对模型中选定的卷积结构执行结构化剪枝，得到初始轻量化模型；

\par\noindent
（3）对剪枝模型进行微调，以恢复其基础检测与分割能力；

\par\noindent
（4）以原始 YOLO11-StackLite 为教师模型，以剪枝模型为学生模型，引入知识蒸馏策略进行进一步训练；

\par\noindent
（5）最终获得兼顾识别性能与模型复杂度的压缩模型。

该流程实现了“先压缩、后补偿”的模型优化思路。其中，结构化剪枝负责降低模型复杂度，知识蒸馏负责缓解精度下降。通过两者结合，可以在尽量保持实例分割与盘点性能的同时，为后续资源受限场景下的模型应用提供更好的部署基础。

\section{本章小结}{Summary}

本章针对 YOLO11-StackLite 在边缘部署需求下仍存在一定参数冗余与计算冗余的问题，提出了结合结构化剪枝与知识蒸馏的模型压缩方案。首先，从边缘场景应用需求出发，分析了模型压缩的必要性；随后，采用结构化剪枝对模型冗余结构进行裁剪，以降低参数量与计算复杂度；在此基础上，引入知识蒸馏策略缓解剪枝带来的性能退化，构建压缩后的轻量化模型。上述方法为后续围绕压缩模型检测、分割与盘点综合性能的实验评价提供了方法基础。


